%%%%%%%%%%%%%%%%%%%%%%%%%%%%%%%%%%%%%%%%%%%%%%%%%%%%%%%%%%%%%%%%%%%%%%%%%%%%%%%%%%
%% SECTION: Opening
%%%%%%%%%%%%%%%%%%%%%%%%%%%%%%%%%%%%%%%%%%%%%%%%%%%%%%%%%%%%%%%%%%%%%%%%%%%%%%%%%%
\section[Overview]{YogaSutra: A Complete Overview}

\begin{frame}[fragile]\frametitle{}
\begin{center}
{\Large YogaSutra: A Complete Overview}\\
{\normalsize पातञ्जल योगसूत्र:Four Chapters, One Journey}
\end{center}
\end{frame}

%%%%%%%%%%%%%%%%%%%%%%%%%%%%%%%%%%%%%%%%%%%%%%%%%%%%%%%%%%%
\begin{frame}[fragile]\frametitle{Why Read the Yoga Sutras?}

\begin{center}
{\large ``If you are on any path where you want to be happy, to be free of suffering, and understand the mind and consciousness, the Sutras are a must-read.''}\\
, \textbf{Prof. Edwin Bryant}, Rutgers University
\end{center}

\vspace{3mm}
\begin{itemize}
\item Written by Sri Patanjali $\approx$ 2000 years ago:the first step-by-step manual for the \emph{practice} of Yoga
\item Most translated ancient Indian text in the medieval era:40+ Indian languages
\item Based on three principles (Nancy Wile):
    \begin{itemize}
    \item Suffering is not caused by events but by our \emph{thoughts about} events
    \item The unlimited peace we seek is our true identity:it exists within us
    \item Peace and Self-realization are attained by \emph{mastering the mind}
    \end{itemize}
\end{itemize}

\tiny{(Ref: Prof. Edwin Bryant; Nancy Wile, Yoga Education Institute; Wikipedia:Yoga Sutras of Patanjali)}

\end{frame}

%%%%%%%%%%%%%%%%%%%%%%%%%%%%%%%%%%%%%%%%%%%%%%%%%%%%%%%%%%%
%% NEW SLIDE A1: What Yoga is NOT
%%%%%%%%%%%%%%%%%%%%%%%%%%%%%%%%%%%%%%%%%%%%%%%%%%%%%%%%%%%
\begin{frame}[fragile]\frametitle{What Yoga is NOT:Clearing the Myths}

\begin{columns}
\begin{column}[t]{0.5\linewidth}
\textbf{Common Misconceptions}
\begin{itemize}
\item ``Yoga = gym stretching'' [No]
\item ``Yoga is a Hindu religion'' [No]
\item ``You must be flexible to do Yoga'' [No]
\item ``It's only for stress relief'' [No]
\item ``Meditation = emptying the mind'' [No]
\end{itemize}
\end{column}
\begin{column}[t]{0.5\linewidth}
\textbf{What Patanjali Actually Says}
\begin{itemize}
\item Yoga = mastery of the \emph{mind}, not the body
\item A universal science of consciousness
\item Available to all:regardless of fitness, religion, or background
\item A complete system: ethics → body → breath → senses → mind → liberation
\item Nirodha ≠ emptying; = \emph{witnessing} without being owned
\end{itemize}
\end{column}
\end{columns}

\vspace{3mm}
\begin{center}
{\small ``The Sutras not only provide a philosophical basis for Yoga, they clarify important esoteric concepts common to all traditions of Indian thought.'':Wikipedia}
\end{center}

\tiny{(Ref: Nancy Wile; Wikipedia:Yoga Sutras of Patanjali; Dr. Ananda Balayogi Bhavanani)}

\end{frame}

%%%%%%%%%%%%%%%%%%%%%%%%%%%%%%%%%%%%%%%%%%%%%%%%%%%%%%%%%%%
%% NEW SLIDE A2: Three Core Principles
%%%%%%%%%%%%%%%%%%%%%%%%%%%%%%%%%%%%%%%%%%%%%%%%%%%%%%%%%%%
\begin{frame}[fragile]\frametitle{Three Core Principles:The Foundation}

\begin{center}
{\normalsize \textbf{Nancy Wile, Yoga Education Institute}}
\end{center}

\vspace{4mm}
\begin{enumerate}
\item \textbf{Suffering is Internal}\\
{\small Suffering is not caused by forces outside us, but by our faulty and limited perception of life and of who we are. Suffering is not caused by the \emph{situation}:but by our \emph{thoughts about} the situation.}

\vspace{3mm}
\item \textbf{Peace is Our True Nature}\\
{\small The unwavering peace we seek is realized by experiencing the unlimited and eternal peace that is our true identity. Though hidden by ignorance, it exists within us, waiting to be revealed.}

\vspace{3mm}
\item \textbf{Mastery of Mind = Self-Realization}\\
{\small Peace and self-realization are attained by mastering the mind. Only a single-pointed, calm mind can reveal the true Self.}
\end{enumerate}

\vspace{2mm}
\begin{center}
{\small These three principles underpin all 196 sutras.}
\end{center}

\tiny{(Ref: Nancy Wile, Yoga Education Institute, 2014)}

\end{frame}

%%%%%%%%%%%%%%%%%%%%%%%%%%%%%%%%%%%%%%%%%%%%%%%%%%%%%%%%%%%
\begin{frame}[fragile]\frametitle{Patanjali \& Structure of YogaSutra}

\begin{columns}
\begin{column}[t]{0.45\linewidth}
\textbf{Maharshi Patanjali पतञ्जलि}
\begin{itemize}
\item ``Father of Yoga'':lived approx. 200–500 BC/AD (scholars differ)
\item Three contributions: Yoga (\textbf{योगसूत्र}), Grammar (\textbf{महाभाष्य}), Medicine (attributed)
\item Did not invent Yoga:\emph{systematized} and \emph{codified} it as सूत्र (aphorisms)
\item Vyasabhasya (commentary) now considered Patanjali's own work
\end{itemize}
\end{column}
\begin{column}[t]{0.55\linewidth}
\textbf{4 पाद (Chapters)}
\begin{center}
{\small
\begin{tabular}{|p{3cm}|r|p{3cm}|}
\hline
\textbf{Chapter} & \textbf{S} & \textbf{Theme} \\
\hline
समाधिपाद Samadhi & 51 & Concentration \\
साधनपाद Sadhana & 55 & Practice \\
विभूतिपाद Vibhuti & 56 & Experiences \\
कैवल्यपाद Kaivalya & 34 & Liberation \\
\hline
\textbf{Total} & \textbf{196} & \\
\hline
\end{tabular}
}
\end{center}
\end{column}
\end{columns}

\vspace{2mm}
\begin{center}
सांख्य is \emph{theory}; Yoga is \emph{practice}:they are the perfect philosophical pair.\\
{\small Samkhya explains \emph{why} we suffer; Yoga shows \emph{how} to be free.}
\end{center}

\tiny{(Ref: Nancy Wile; Dr. Ananda Balayogi Bhavanani; Wikipedia)}

\end{frame}

%%%%%%%%%%%%%%%%%%%%%%%%%%%%%%%%%%%%%%%%%%%%%%%%%%%%%%%%%%%
\begin{frame}[fragile]\frametitle{Key Yogic Concepts}
\begin{itemize}
\item \textbf{पुरुष (Purusha)}:Pure consciousness; the Seer; eternal, unchanging witness
\item \textbf{प्रकृति (Prakriti)}:Primordial nature; composed of three गुण (सत्व, रजस्, तमस्)
\item \textbf{चित्त (Chitta)}:Mind-stuff; four coordinated functions:
    \begin{itemize}
    \item \textbf{मनस्}:sensory receiver; processes perception and directs action
    \item \textbf{बुद्धि}:intellect; faculty of discrimination and decision
    \item \textbf{अहंकार}:ego; the ``I-maker'' that claims ``I think, I act''
    \item \textbf{चित्त} (memory store):repository of संस्कार (impressions)
    \end{itemize}
\item \textbf{वृत्ति (Vritti)}:modifications/ripples of mind; the disturbance Yoga addresses
\item \textbf{क्लेश (Klesha)}:deep afflictions that generate vrittis and perpetuate suffering
\item \textbf{कैवल्य (Kaivalya)}:liberation; Purusha abiding in its own pure nature
\end{itemize}

\vspace{2mm}
\begin{center}
{\small Chitta is like a river, flowing between worldliness and Kaivalya. (Vyasa)\\
``We learn to control our mind and not be controlled by our thoughts.'':Nancy Wile}
\end{center}

\tiny{(Ref: Vyasa Bhashya; Krishnamurthy Ramakrishnan; Dr. Veena Londhe)}

\end{frame}

%%%%%%%%%%%%%%%%%%%%%%%%%%%%%%%%%%%%%%%%%%%%%%%%%%%%%%%%%%%%%%%%%%%%%%%%%%%%%%%%%%
%% SECTION: Chapter 1:Samadhi Pada
%%%%%%%%%%%%%%%%%%%%%%%%%%%%%%%%%%%%%%%%%%%%%%%%%%%%%%%%%%%%%%%%%%%%%%%%%%%%%%%%%%
\section[Samadhipad]{Chapter 1: Samadhi Pada:Concentration}

\begin{frame}[fragile]\frametitle{}
\begin{center}
{\Large Chapter 1: समाधिपाद Samadhi Pada}\\
{\normalsize Concentration:51 Sutras}\\
{\small For \textbf{uttama} (superior) seekers already near the goal; addresses the nature of mind, the goal of Yoga, and the twin keys of practice and non-attachment.}\\
{\small The first 16 sutras (1.1–1.16) are sometimes called the \textbf{``Sweet Sixteen''}:the complete philosophy of Yoga in a nutshell.}
\end{center}
\tiny{(Ref: Nancy Wile, Yoga Education Institute; Dr. Ananda Balayogi Bhavanani)}
\end{frame}

%%%%%%%%%%%%%%%%%%%%%%%%%%%%%%%%%%%%%%%%%%%%%%%%%%%%%%%%%%%
\begin{frame}[fragile]\frametitle{1.1:अथ योगानुशासनम्}

\begin{sanskrit}
अथ योगानुशासनम् ॥१.१॥
\end{sanskrit}

\begin{itemize}
\item \textbf{अथ}:not a casual ``now''; signals \emph{existential readiness}: ethical maturity, intellectual clarity, earnest longing for liberation
    \begin{itemize}
    \item ``Now, after having done prior preparation through life and other practices, the study and practice of Yoga begins.'' (Nancy Wile)
    \item You are coming from different evolutionary paths; \emph{now} you are eligible
    \end{itemize}
\item \textbf{अनुशासनम्}:systematic transmission; a living tradition being re-organized, not invented; \emph{अनु} = following the existing tradition
\item \textbf{Key insight:} Yoga does not begin with technique or posture:it begins the moment one becomes \emph{ready to be taught}
\end{itemize}

\tiny{(Ref: Krishnamurthy Ramakrishnan; Nancy Wile; Vyasa Bhashya; Prof. Edwin Bryant)}

\end{frame}

%%%%%%%%%%%%%%%%%%%%%%%%%%%%%%%%%%%%%%%%%%%%%%%%%%%%%%%%%%%
\begin{frame}[fragile]\frametitle{1.2:योगश्चित्तवृत्तिनिरोधः}

\begin{sanskrit}
योगश् चित्तवृत्तिनिरोधः ॥१.२॥
\end{sanskrit}

\begin{itemize}
\item The definition of Yoga in four words:not postures, not breathing, but \emph{stilling of mind-modifications}
\item \textbf{चित्त}: four coordinated voices:मनस् receives → बुद्धि evaluates → चित्त supplies memory → अहंकार asserts ``I am doing this''
\item \textbf{वृत्ति}: ripples on the lake of the mind:the disturbances, not the lake itself
\item \textbf{निरोधः}: restraint, regulation, coordination:not suppression (उच्छेद); root \emph{रुध्} = to check movement
\item \textbf{5 states of Chitta:} मूढ (dull) → क्षिप्त (restless) → विक्षिप्त (partly steady) → एकाग्र (focused) → निरुद्ध (fully restrained)
\item Edwin Bryant: ``Yoga is meditative practices culminating in a state of consciousness free from all modes of active or discursive thought.''
\end{itemize}

\tiny{(Ref: Krishnamurthy Ramakrishnan; [SS] Swami Satchidananda; Prof. Edwin Bryant; Vyasa Bhashya)}

\end{frame}

%%%%%%%%%%%%%%%%%%%%%%%%%%%%%%%%%%%%%%%%%%%%%%%%%%%%%%%%%%%
\begin{frame}[fragile]\frametitle{1.3–1.4:The Seer and the Seen}

\begin{sanskrit}
तदा द्रष्टुः स्वरूपेऽवस्थानम् ॥१.३॥
वृत्तिसारूप्यमितरत्र ॥१.४॥
\end{sanskrit}

\textbf{1.3}: ``Then the Seer abides in Itself, resting in its own True Nature, which is called Self-realization.'' (Nancy Wile)\\
\textbf{1.4}:At all other times, the Seer appears to take the form of the vrittis:it takes on the identity of those thought patterns.

\begin{itemize}
\item \textbf{Core metaphor:} Mind = lake; Purusha = the still bottom; Vrittis = ripples. As long as the lake is disturbed, one cannot see the bottom clearly.
\item Yoga stills the lake so the true depth (Purusha) is revealed.
\item \textbf{Identity question:} ``What identities do you give yourself?'':your job, your role, your story. Are those really \emph{you}? (Nancy Wile)
\item We do not \emph{create} the Seer:we \emph{uncover} what was always there.
\end{itemize}

\tiny{(Ref: Nancy Wile; [VH] Vyasa Houston; [BM] Barbara Miller; Swami Vivekananda)}

\end{frame}

%%%%%%%%%%%%%%%%%%%%%%%%%%%%%%%%%%%%%%%%%%%%%%%%%%%%%%%%%%%
\begin{frame}[fragile]\frametitle{1.5–1.6:The Five Vrittis}

\begin{sanskrit}
वृत्तयः पञ्चतय्यः क्लिष्टाक्लिष्टाः ॥१.५॥
प्रमाणविपर्ययविकल्पनिद्रास्मृतयः ॥१.६॥
\end{sanskrit}

\textbf{Five types of modifications:all require restraint:}

\begin{center}
{\small
\begin{tabular}{|p{2cm}|p{2.5cm}|p{2.5cm}|}
\hline
\textbf{Vritti} & \textbf{Meaning} & \textbf{Example} \\
\hline
प्रमाण & Valid knowledge & Direct seeing \\
विपर्यय & Misperception & Rope seen as snake \\
विकल्प & Imagination & ``Rabbit's horn'' \\
निद्रा & Sleep & Absence-based vritti \\
स्मृति & Memory & Retention of impressions \\
\hline
\end{tabular}
}
\end{center}

\begin{itemize}
\item \textbf{क्लिष्ट} (pain-afflicted) vs. \textbf{अक्लिष्ट} (not pain-afflicted): not a moral division; a phenomenological tone
\item Even pleasant vrittis disturb equilibrium:all require restraint
\end{itemize}

\tiny{(Ref: Krishnamurthy Ramakrishnan; Nancy Wile; [IT] I.K. Taimni)}

\end{frame}

%%%%%%%%%%%%%%%%%%%%%%%%%%%%%%%%%%%%%%%%%%%%%%%%%%%%%%%%%%%
%% NEW SLIDE: How a Vritti Forms
%%%%%%%%%%%%%%%%%%%%%%%%%%%%%%%%%%%%%%%%%%%%%%%%%%%%%%%%%%%
\begin{frame}[fragile]\frametitle{Before the Ripple Surfaces:How a Vritti (वृत्ति) Forms}

\begin{itemize}
\item Chitta (चित्त) holds far more activity than we are ever aware of:most of it stays below the threshold of noticing
\item A Vritti (वृत्ति) begins as an impulse stirring in that depth:not yet a ``thought'' in the way we experience thought
\item It becomes a perceived thought only when attention falls on it:मनस् (Manas) receiving + बुद्धि (Buddhi) evaluating together illuminate a few impulses out of many
\item This is why two people in the same situation can have entirely different thoughts:different impulses get lit up
\item The rising happens largely on its own, prior to conscious choice:we do not decide which impulse surfaces
\end{itemize}

\tiny{(Ref: Swami Vivekananda's lake metaphor; Vritti:Wikipedia; Yoga Sutra 1.2 commentary tradition)}

\end{frame}

%%%%%%%%%%%%%%%%%%%%%%%%%%%%%%%%%%%%%%%%%%%%%%%%%%%%%%%%%%%
%% NEW SLIDE: Why Nirodha Targets the Rising
%%%%%%%%%%%%%%%%%%%%%%%%%%%%%%%%%%%%%%%%%%%%%%%%%%%%%%%%%%%
\begin{frame}[fragile]\frametitle{Why Nirodha (निरोध) Targets the Rising, Not the Fight}

\begin{itemize}
\item Since a Vritti (वृत्ति) surfaces before we consciously choose it, fighting it \emph{after} it has risen is already fighting an effect, not the cause
\item Nirodha (निरोध) works upstream:training attention so fewer disturbing impulses get illuminated in the first place
\item Not a suppression of activity (उच्छेद, uchchheda), but a steady redirection:root रुध् (rudh) = to check movement, not to destroy it
\item Memory aid only, not a scriptural distinction:opposing an outside circumstance (loosely, विरोध, virodha) is a different move from this inward, self-directed restraint (निरोध, nirodha)
\item Restraint is carried out through Abhyasa (अभ्यास) and Vairagya (वैराग्य)
\end{itemize}

\tiny{(Ref: root रुध् etymology; Yoga Sutra 1.2 commentary tradition; forward-reference to 1.12)}

\end{frame}

%%%%%%%%%%%%%%%%%%%%%%%%%%%%%%%%%%%%%%%%%%%%%%%%%%%%%%%%%%%
%% NEW SLIDE B1: Colored vs Uncolored Thoughts
%%%%%%%%%%%%%%%%%%%%%%%%%%%%%%%%%%%%%%%%%%%%%%%%%%%%%%%%%%%
\begin{frame}[fragile]\frametitle{``Coloring'' Your Thoughts:1.5 in Daily Life}

\begin{center}
{\normalsize \textbf{Many of our thoughts are colored by our judgments and expectations.}}
\end{center}

\begin{itemize}
\item A car is driving down the street = \textbf{fact} (प्रमाण)
\item That car is mine / I want it = \textbf{concept} (colored by desire = राग)
\item Seeing a squirrel run in front of the car, you brake:but it's just dry leaves (विपर्यय)
\end{itemize}

\textbf{Practical practice from Nancy Wile:}
\begin{itemize}
\item When a thought arises, simply note: \emph{``This is colored''} or \emph{``This is not colored''}
\item Ask: \emph{``Is this useful or not useful?''}:observe without fighting or feeding the thought
\item ``These sutras focus on not fighting against anxiety or depression, but just \emph{letting it be} and observing it.'':Nancy Wile
\item Even physical pain: suffering reduces when we observe it with detachment rather than adding a \emph{story} around the pain
\end{itemize}

\begin{center}
{\small The mind may want to make something a problem. It's best to just observe.}
\end{center}

\tiny{(Ref: Nancy Wile, Yoga Education Institute, p.6-7)}

\end{frame}

%%%%%%%%%%%%%%%%%%%%%%%%%%%%%%%%%%%%%%%%%%%%%%%%%%%%%%%%%%%
\begin{frame}[fragile]\frametitle{1.7–1.9:Valid Knowledge, Misperception, Imagination}

\begin{sanskrit}
प्रत्यक्षानुमानागमाः प्रमाणानि ॥१.७॥
विपर्ययो मिथ्याज्ञानमतद्रूपप्रतिष्ठम् ॥१.८॥
शब्दज्ञानानुपाती वस्तुशून्यो विकल्पः ॥१.९॥
\end{sanskrit}

\begin{itemize}
\item \textbf{1.7:Three Pramanas (valid knowledge):}
    \begin{itemize}
    \item \textbf{प्रत्यक्ष}:direct perception
    \item \textbf{अनुमान}:inference (fire from smoke)
    \item \textbf{आगम}:authoritative testimony (reliable teacher or shastra)
    \end{itemize}
\item \textbf{1.8:विपर्यय (Misperception):} false knowledge not grounded in the object's true nature; rope → snake; emotional reactions arise from misperception, not reality
\item \textbf{1.9:विकल्प (Imagination):} word-based cognition without a real object; ``horn of a rabbit''; modern social media massively amplifies vikalpa:we react to images of things that never happened to us
\end{itemize}

\tiny{(Ref: Krishnamurthy Ramakrishnan; Nancy Wile; Patanjaliyogasutra.in)}

\end{frame}

%%%%%%%%%%%%%%%%%%%%%%%%%%%%%%%%%%%%%%%%%%%%%%%%%%%%%%%%%%%
\begin{frame}[fragile]\frametitle{1.10–1.11:Sleep and Memory}

\begin{sanskrit}
अभावप्रत्ययालम्बना वृत्तिर्निद्रा ॥१.१०॥
अनुभूतविषयासंप्रमोषः स्मृतिः ॥१.११॥
\end{sanskrit}

\textbf{1.10:निद्रा (Sleep):}
\begin{itemize}
\item Sleep is not the absence of mind:it is a vritti whose object is \emph{non-existence}
\item ``I slept well'' proves awareness persisted in seed form
\item Deep sleep rehearses Samadhi:the difference: Samadhi brings \emph{awareness} into that stillness
\end{itemize}

\textbf{1.11:स्मृति (Memory):}
\begin{itemize}
\item Memory = \emph{non-loss} of experienced objects; faithful retention of संस्कार (impressions)
\item Tool: enables learning; Trap: fixation on past pain
\item ``We create our own world through our mental projections.'':Nancy Wile
\item Liberation lies \emph{beyond} both memory and forgetting
\end{itemize}

\tiny{(Ref: Krishnamurthy Ramakrishnan; Nancy Wile; [HA] Hariharananda Aranya; [VH] Vyasa Houston)}

\end{frame}

%%%%%%%%%%%%%%%%%%%%%%%%%%%%%%%%%%%%%%%%%%%%%%%%%%%%%%%%%%%
\begin{frame}[fragile]\frametitle{1.12:अभ्यासवैराग्याभ्यां तन्निरोधः}

\begin{sanskrit}
अभ्यासवैराग्याभ्यां तन्निरोधः ॥१.१२॥
\end{sanskrit}

\begin{center}
{\large Restraint of vrittis is achieved through \textbf{Practice} and \textbf{Non-Attachment}}
\end{center}

\begin{columns}
\begin{column}[t]{0.5\linewidth}
\textbf{अभ्यास (Abhyasa):Practice}
\begin{itemize}
\item Choosing and applying effort consistently
\item Brings \emph{steadiness} (sthiti) to the mind
\item The active, disciplined wing
\end{itemize}
\end{column}
\begin{column}[t]{0.5\linewidth}
\textbf{वैराग्य (Vairagya):Non-Attachment}
\begin{itemize}
\item Cultivated dispassion:not indifference
\item Letting go of attachment to results and experiences
\item The releasing, freeing wing
\end{itemize}
\end{column}
\end{columns}

\vspace{3mm}
\begin{center}
{\small Like the two wings of a bird:both are needed for flight.\\
\emph{``Practice and non-attachment are the twin keys upon which the entire system of Yoga rests.''}:Dr. Ananda}
\end{center}

\tiny{(Ref: Dr. Ananda Balayogi Bhavanani; [SS] Swami Satchidananda; Krishnamurthy Ramakrishnan)}

\end{frame}

%%%%%%%%%%%%%%%%%%%%%%%%%%%%%%%%%%%%%%%%%%%%%%%%%%%%%%%%%%%
\begin{frame}[fragile]\frametitle{1.13–1.14:Defining Practice and Making it Firm}

\begin{sanskrit}
तत्र स्थितौ यत्नोऽभ्यासः ॥१.१३॥
स तु दीर्घकालनैरन्तर्यसत्कारासेवितो दृढभूमिः ॥१.१४॥
\end{sanskrit}

\textbf{1.13:Abhyasa defined:} the effort directed toward steadiness (sthiti); choosing actions that bring a stable and tranquil state.

\textbf{1.14:Practice becomes a firm foundation when done with:}
\begin{itemize}
\item \textbf{दीर्घकाल}:for a \emph{long time}; not weeks, but years of sustained commitment
\item \textbf{नैरन्तर्य}:\emph{without break}; continuity; Yoga practice should be steady and without gaps
\item \textbf{सत्कार}:with \emph{earnestness and reverence}; sincere devotion, not mechanical performance
\item \textbf{आसेवित}:faithfully attended to; the whole heart behind the practice
\end{itemize}

\vspace{2mm}
\begin{center}
{\small ``Yoga practice should be steady and without gaps.'':Nancy Wile (1.14)\\
``Yoga is not a shortcut. It is a long road walked with joy.''}
\end{center}

\tiny{(Ref: Nancy Wile; [HA] Hariharananda Aranya; Patanjaliyogasutra.in)}

\end{frame}

%%%%%%%%%%%%%%%%%%%%%%%%%%%%%%%%%%%%%%%%%%%%%%%%%%%%%%%%%%%
\begin{frame}[fragile]\frametitle{1.15–1.16:Vairagya and Para-Vairagya}

\begin{sanskrit}
दृष्टानुश्रविकविषयवितृष्णस्य वशीकारसंज्ञा वैराग्यम् ॥१.१५॥
तत्परं पुरुषख्यातेर्गुणवैतृष्ण्यम् ॥१.१६॥
\end{sanskrit}

\textbf{1.15:Vairagya (Non-attachment):}
\begin{itemize}
\item The mind that has lost craving for objects seen \emph{or} described has achieved vairagya
\item Not hatred of life:cultivated \emph{dispassionate objectivity}
\item ``We give up attachments to things, then to experiences, then to people, then to our own identity.'' (Nancy Wile)
\item Many people identify with labels:job, role, place. Notice what happens when you let go of your story. If they go away, you are still here.
\end{itemize}

\textbf{1.16:Para-Vairagya (Supreme Non-attachment):}
\begin{itemize}
\item Indifference even to the three \textbf{gunas}, achieved through knowledge of \textbf{Purusha}
\item Letting go of even subtle spiritual pleasures and yogic attainments
\item Destroys the very \emph{seeds} of the impurities
\end{itemize}

\tiny{(Ref: Dr. Ananda Balayogi Bhavanani; Nancy Wile; [IT] I.K. Taimni)}

\end{frame}

%%%%%%%%%%%%%%%%%%%%%%%%%%%%%%%%%%%%%%%%%%%%%%%%%%%%%%%%%%%
%% NEW SLIDE B2: 1.17-1.22 Five Powers + Seeker Types
%%%%%%%%%%%%%%%%%%%%%%%%%%%%%%%%%%%%%%%%%%%%%%%%%%%%%%%%%%%
\begin{frame}[fragile]\frametitle{1.17–1.22:Samadhi Types, Five Powers \& Seeker Intensity}

\textbf{1.17–1.18:Types of Samadhi (Samprajnata / Asamprajnata):}
\begin{itemize}
\item Four levels of object-based absorption: \textbf{Vitarka} (gross) → \textbf{Vichara} (subtle) → \textbf{Ananda} (bliss) → \textbf{Asmita} (I-am-ness):like listening to a song: you first notice it, then examine it, then feel it, then lose yourself in it (Nancy Wile)
\item Asamprajnata (1.18): objectless; only latent impressions remain
\end{itemize}

\textbf{1.20:Five Powers (Panchabala) for the systematic seeker:}
\begin{enumerate}
\item \textbf{श्रद्धा} Shraddha:faith in one's own power to reach enlightenment
\item \textbf{वीर्य} Veerya:strength of body and mind; vigour
\item \textbf{स्मृति} Smriti:repeated memory of the path; keeping the goal alive
\item \textbf{समाधि} Samadhi-prajna:training in deep concentration
\item \textbf{प्रज्ञा} Prajna:pursuit of real knowledge; wisdom
\end{enumerate}

\textbf{1.21–1.22:Intensity matters:}
\begin{itemize}
\item Mild / Medium / Intense practitioners reach the goal at different speeds
\item ``Those who pursue with intensity of feeling, vigour, and firm conviction achieve concentration more quickly.'' (1.21)
\end{itemize}

\tiny{(Ref: Nancy Wile p.8-9; [HA] Hariharananda Aranya; Dr. Ananda Balayogi Bhavanani)}

\end{frame}

%%%%%%%%%%%%%%%%%%%%%%%%%%%%%%%%%%%%%%%%%%%%%%%%%%%%%%%%%%%
\begin{frame}[fragile]\frametitle{1.23 \& 1.28–1.29:The Direct Route: Ishvara and OM}

\begin{sanskrit}
ईश्वरप्रणिधानाद्वा ॥१.२३॥
तज्जपस्तदर्थभावनम् ॥१.२८॥
\end{sanskrit}

\textbf{1.23:Ishvara Pranidhana:}
\begin{itemize}
\item Or, through \emph{dedication and surrender to Ishvara}:the direct route for devotional seekers
\item Ishvara = a special Purusha untouched by klesha, karma, and time; the eternal teacher (1.24–1.26)
\item ``Ishvara is anything that has spiritual significance to the individual'':non-sectarian, universal (Wikipedia)
\item For those whose heart leads the way, this is the fastest path
\end{itemize}

\textbf{1.27–1.29:Pranava (OM):}
\begin{itemize}
\item The sacred word (प्रणव) designating Ishvara is the sound OM
\item OM = three sounds: \textbf{अ} (creation) + \textbf{उ} (sustenance) + \textbf{म} (dissolution)
\item Its repetition (\textbf{जप}) and meditation on its meaning brings realization of the individual Self and removal of obstacles (1.29)
\end{itemize}

\tiny{(Ref: Dr. Ananda Balayogi Bhavanani; Nancy Wile; Wikipedia; [SS] Swami Satchidananda)}

\end{frame}

%%%%%%%%%%%%%%%%%%%%%%%%%%%%%%%%%%%%%%%%%%%%%%%%%%%%%%%%%%%
\begin{frame}[fragile]\frametitle{1.30–1.32:The Nine Obstacles}

\begin{sanskrit}
व्याधिस्त्यानसंशयप्रमादालस्याविरतिभ्रान्तिदर्शनालब्धभूमिकत्वानवस्थितत्वानि चित्तविक्षेपाः तेऽन्तरायाः ॥१.३०॥
\end{sanskrit}

\textbf{Nine obstacles (अन्तराय):predictable, not personal failures:}

\begin{enumerate}
\item व्याधि:disease \quad 2. स्त्यान:dullness \quad 3. संशय:doubt
\item प्रमाद:carelessness \quad 5. आलस्य:laziness \quad 6. अविरति:sensuality
\item भ्रान्तिदर्शन:false perception \quad 8. अलब्धभूमिकत्व:failure to reach a stage
\item अनवस्थितत्व:inability to maintain a stage
\end{enumerate}

\textbf{1.31:Four companions (symptoms):} Sorrow, mental agitation, body trembling, disturbed breathing\\
\textbf{1.32:One remedy:} Single-pointed practice of \emph{one principle} is the antidote for all nine. ``Doubt, including worrying about what others think and worrying about the future, can be a major obstacle.'':Nancy Wile

\tiny{(Ref: Nancy Wile p.11; Krishnamurthy Ramakrishnan; Dr. Ananda Balayogi Bhavanani)}

\end{frame}

%%%%%%%%%%%%%%%%%%%%%%%%%%%%%%%%%%%%%%%%%%%%%%%%%%%%%%%%%%%
\begin{frame}[fragile]\frametitle{1.33:Four Attitudes for a Clear Mind}

\begin{sanskrit}
मैत्रीकरुणामुदितोपेक्षाणां सुखदुःखपुण्यापुण्यविषयाणां भावनातश्चित्तप्रसादनम् ॥१.३३॥
\end{sanskrit}

\begin{itemize}
\item \textbf{मैत्री} (Maitri):friendliness toward the happy
\item \textbf{करुणा} (Karuna):compassion toward those suffering
\item \textbf{मुदिता} (Mudita):joy toward the virtuous
\item \textbf{उपेक्षा} (Upeksha):equanimity toward the non-virtuous
\end{itemize}

\vspace{2mm}
\begin{itemize}
\item \textbf{Why this works:} By serving others and thinking about their needs, we eliminate our own unending wanting. ``Once we are free from attachment to our own desires, we can find true and lasting happiness.'':Nancy Wile
\item This sutra appears in Buddhist texts too:a universal human wisdom (Brahmaviharas)
\item \textbf{1.34:Pranayama:} The mind is also calmed by regulating the breath, attending especially to the exhalation
\end{itemize}

\tiny{(Ref: Nancy Wile p.12; Dr. Ananda Balayogi Bhavanani; Wikipedia:Brahmaviharas)}

\end{frame}

%%%%%%%%%%%%%%%%%%%%%%%%%%%%%%%%%%%%%%%%%%%%%%%%%%%%%%%%%%%
%% NEW SLIDE B3: 1.34-1.39 Six Paths
%%%%%%%%%%%%%%%%%%%%%%%%%%%%%%%%%%%%%%%%%%%%%%%%%%%%%%%%%%%
\begin{frame}[fragile]\frametitle{1.34–1.39:Six Paths to Stabilize the Mind}

\begin{center}
{\normalsize \textbf{Choose the method that resonates with you:Patanjali offers six:}}
\end{center}

\begin{center}
{\small
\begin{tabular}{|p{0.8cm}|p{3cm}|p{3.2cm}|}
\hline
\textbf{S} & \textbf{Method} & \textbf{Modern parallel} \\
\hline
1.34 & Pranayama: regulate the breath & Box breathing, 4-7-8 technique \\
1.35 & Inner concentration on sensory process & Body-scan meditation \\
1.36 & Concentration on inner light / luminosity & Trataka, inner light practice \\
1.37 & Contemplating a desireless mind & Witnessing without craving \\
1.38 & Focusing on dream state or deep sleep & Yoga Nidra, lucid dreaming \\
1.39 & Whatever object or principle you love & Devotional focus on your chosen ideal \\
\hline
\end{tabular}
}
\end{center}

\vspace{2mm}
\begin{center}
{\small ``Or by concentrating on \emph{whatever object} one may like, or towards which one has a predisposition, the mind becomes stable and tranquil.'' (1.39, Nancy Wile)\\
\textbf{Yoga meets you where you are.}}
\end{center}

\tiny{(Ref: Nancy Wile p.12; [SS] Swami Satchidananda; Dr. Ananda Balayogi Bhavanani)}

\end{frame}

%%%%%%%%%%%%%%%%%%%%%%%%%%%%%%%%%%%%%%%%%%%%%%%%%%%%%%%%%%%
%% NEW SLIDE B4: 1.41 Crystal Mind
%%%%%%%%%%%%%%%%%%%%%%%%%%%%%%%%%%%%%%%%%%%%%%%%%%%%%%%%%%%
\begin{frame}[fragile]\frametitle{1.41:The Crystal Mind \& Stages of Samadhi (1.40–1.51)}

\begin{sanskrit}
क्षीणवृत्तेरभिजातस्येव मणेर्ग्रहीतृग्रहणग्राह्येषु तत्स्थतदञ्जनता समापत्तिः ॥१.४१॥
\end{sanskrit}

\begin{itemize}
\item \textbf{1.40:} When vrittis are mastered, the mind can hold any object:from the smallest atom to the infinite
\item \textbf{1.41:Sphatika (Crystal) Mind:} When vrittis weaken, the mind becomes like a \textbf{transparent crystal}:it takes on the exact quality of whatever it rests on; observer, observing, and observed merge into one (समापत्ति)
\item Modern psychology calls this \emph{metacognition} or \emph{witness consciousness}:the ability to observe thought without being swept away by it
\end{itemize}

\textbf{Samadhi:four stages (1.42–1.51):}
\begin{itemize}
\item \textbf{सम्प्रज्ञात} (Samprajnata): object-based:Vitarka (gross) → Vichara (subtle) → Ananda (bliss) → Asmita (I-am-ness)
\item \textbf{असम्प्रज्ञात} (Asamprajnata): objectless; only latent impressions remain
\item \textbf{निर्बीज} (Nirbija, 1.51): seedless; all vrittis and samskaras dissolved
\end{itemize}

\tiny{(Ref: Nancy Wile p.13; Dr. Ananda Balayogi Bhavanani; [IT] I.K. Taimni)}

\end{frame}

%%%%%%%%%%%%%%%%%%%%%%%%%%%%%%%%%%%%%%%%%%%%%%%%%%%%%%%%%%%%%%%%%%%%%%%%%%%%%%%%%%
%% SECTION: Chapter 2:Sadhana Pada
%%%%%%%%%%%%%%%%%%%%%%%%%%%%%%%%%%%%%%%%%%%%%%%%%%%%%%%%%%%%%%%%%%%%%%%%%%%%%%%%%%
\section[Sadhanapad]{Chapter 2: Sadhana Pada:Practice}

\begin{frame}[fragile]\frametitle{Chapter 2: साधनपाद Sadhana Pada:Practice (55 Sutras)}

\begin{itemize}
\item \textbf{Theme:} For \emph{madhyama/adham} (ordinary) seekers; practical tools to reduce suffering and prepare the ground for Samadhi
\item Two yogic paths:
    \begin{enumerate}
    \item \textbf{क्रियायोग (Kriya Yoga)} (2.1):Yoga of Action: cleanses deeply rooted afflictions through tapas, svadhyaya, and Ishvara-pranidhana
    \item \textbf{अष्टांगयोग (Ashtanga Yoga)} (2.28–2.29):8-limbed path: systematic outer-to-inner progression
    \end{enumerate}
\item The chapter also gives Yoga's most important diagnosis: the \textbf{five kleshas} (2.3):root causes of all suffering
\item ``Pada Two focuses on motivation for regular practice while offering clear, comprehensive instructions.'':Nancy Wile
\end{itemize}

\tiny{(Ref: Nancy Wile; Dr. Ananda Balayogi Bhavanani; Wikipedia)}

\end{frame}

%%%%%%%%%%%%%%%%%%%%%%%%%%%%%%%%%%%%%%%%%%%%%%%%%%%%%%%%%%%
%% NEW SLIDE C1: The Three Gunas
%%%%%%%%%%%%%%%%%%%%%%%%%%%%%%%%%%%%%%%%%%%%%%%%%%%%%%%%%%%
\begin{frame}[fragile]\frametitle{The Three Gunas:Modes of Mind and Nature}

\begin{center}
{\normalsize Prakriti (nature) and all mental states are composed of three \textbf{गुण} (qualities):}
\end{center}

\begin{center}
{\small
\begin{tabular}{|p{1.5cm}|p{1.5cm}|p{1.5cm}|p{2.5cm}|}
\hline
\textbf{Guna} & \textbf{Quality} & \textbf{Mind-state} & \textbf{Daily experience} \\
\hline
\textbf{सत्व} Sattva & Clarity, harmony & Alert, focused, joyful & Creative flow, good sleep, insight \\
\textbf{रजस्} Rajas & Passion, activity & Restless, craving, driven & Anxiety, FOMO, overwork, anger \\
\textbf{तमस्} Tamas & Inertia, dullness & Sluggish, confused & Procrastination, depression, lethargy \\
\hline
\end{tabular}
}
\end{center}

\vspace{2mm}
\begin{itemize}
\item All three are always present; our mental state depends on which predominates
\item Yoga practice gradually increases Sattva until even Sattva is transcended (Para-vairagya, 1.16)
\item \textbf{Recognition test:} Which guna describes your mind right now? Before sleep? After heavy food? After meditation?
\end{itemize}

\begin{center}
{\small ``When Sattva predominates: lucidity, wisdom, harmony, peacefulness. When Rajas: craving, restlessness. When Tamas: ignorance, lethargy, suffering.'':Wikipedia}
\end{center}

\tiny{(Ref: Wikipedia:Yoga Sutras of Patanjali; Dr. Ananda Balayogi Bhavanani; [IT] I.K. Taimni)}

\end{frame}

%%%%%%%%%%%%%%%%%%%%%%%%%%%%%%%%%%%%%%%%%%%%%%%%%%%%%%%%%%%
\begin{frame}[fragile]\frametitle{2.1–2.2:Kriya Yoga: The Yoga of Action}

\begin{sanskrit}
तपःस्वाध्यायेश्वरप्रणिधानानि क्रियायोगः ॥२.१॥
समाधिभावनार्थः क्लेशतनूकरणार्थश्च ॥२.२॥
\end{sanskrit}

\textbf{2.1:Three limbs of Kriya Yoga:}
\begin{itemize}
\item \textbf{तपस् (Tapas)}:Endurance; acceptance of hardship as purification; ``accepting pain as help for purification'' [SS]; training and purifying the senses
\item \textbf{स्वाध्याय (Svadhyaya)}:Self-study; study of sacred texts; honest inner observation; going to the root cause of one's patterns; self-study in the context of teachings
\item \textbf{ईश्वरप्रणिधान (Ishvara Pranidhana)}:Surrender to the Divine; offering the fruits of all actions; devotion as a daily practice; letting go into the creative source from which we emerged
\end{itemize}

\textbf{2.2:Purpose:}
\begin{itemize}
\item To attain \textbf{Samadhi} and to weaken and dissolve the \textbf{Kleshas}
\end{itemize}

\tiny{(Ref: [HA] Hariharananda Aranya; [SS] Swami Satchidananda; Nancy Wile; Dr. Ananda Balayogi Bhavanani)}

\end{frame}

%%%%%%%%%%%%%%%%%%%%%%%%%%%%%%%%%%%%%%%%%%%%%%%%%%%%%%%%%%%
\begin{frame}[fragile]\frametitle{2.3–2.5:The Five Kleshas: Root Causes of Suffering}

\begin{sanskrit}
अविद्यास्मितारागद्वेषाभिनिवेशाः क्लेशाः ॥२.३॥
अविद्याक्षेत्रमुत्तरेषाम् ॥२.४॥
\end{sanskrit}

\textbf{Five Afflictions (क्लेश):Patanjali's root-cause diagnosis:}

\begin{center}
{\small
\begin{tabular}{|p{2cm}|p{2cm}|p{3cm}|}
\hline
\textbf{Klesha} & \textbf{Meaning} & \textbf{Daily life form} \\
\hline
\textbf{अविद्या} Avidya & Ignorance of true Self & Mistaking the body/mind for the self \\
\textbf{अस्मिता} Asmita & Egoism, I-ness & Over-identification with roles \\
\textbf{राग} Raga & Attachment & Craving for pleasant experiences \\
\textbf{द्वेष} Dvesha & Aversion & Repulsion of unpleasant experiences \\
\textbf{अभिनिवेश} Abhinivesha & Fear of death & Clinging to life and what we know \\
\hline
\end{tabular}
}
\end{center}

\begin{itemize}
\item \textbf{2.4:} Avidya is the root-field of all other four:remove the root and all branches wither
\item \textbf{2.5:} Avidya = four confusions: transient as eternal; impure as pure; painful as pleasant; not-self as Self
\end{itemize}

\tiny{(Ref: [IT] I.K. Taimni; Dr. Ananda Balayogi Bhavanani; Nancy Wile)}

\end{frame}

%%%%%%%%%%%%%%%%%%%%%%%%%%%%%%%%%%%%%%%%%%%%%%%%%%%%%%%%%%%
%% NEW SLIDE C2: 2.16 + Pratipaksha Bhavana
%%%%%%%%%%%%%%%%%%%%%%%%%%%%%%%%%%%%%%%%%%%%%%%%%%%%%%%%%%%
\begin{frame}[fragile]\frametitle{2.16 \& 2.33:Prevention \& the Counter-Thought Practice}

\begin{sanskrit}
हेयं दुःखमनागतम् ॥२.१६॥
वितर्कबाधने प्रतिपक्षभावनम् ॥२.३३॥
\end{sanskrit}

\textbf{2.16:Yoga's Preventive Medicine:}
\begin{itemize}
\item ``The pain \emph{yet to come} can be avoided.'':one of the most action-oriented sutras
\item Patanjali advises us to prevent miseries \emph{before} they arise
\item ``It is the pain which is yet to come that is to be avoided and discarded.'' (Nancy Wile)
\item How we respond to events now determines our future experience
\end{itemize}

\textbf{2.33:Pratipaksha Bhavana (Counter-thought):}
\begin{itemize}
\item When negative thoughts cause suffering, deliberately cultivate the \emph{opposite}
\item Even if we cannot replace a thought with positive emotion, at minimum make a strong \textbf{STOP}:this works wonders in blocking negative spirals (Dr. Ananda)
\item \textbf{Practical:} Anxiety → deliberately recall a moment of genuine calm. Anger → deliberately imagine the other person's perspective.
\item Modern parallel: Cognitive Behavioural Therapy (CBT) reframing technique
\end{itemize}

\tiny{(Ref: Dr. Ananda Balayogi Bhavanani; Nancy Wile p.16; [SS] Swami Satchidananda)}

\end{frame}

%%%%%%%%%%%%%%%%%%%%%%%%%%%%%%%%%%%%%%%%%%%%%%%%%%%%%%%%%%%
%% NEW SLIDE C3: Yamas & Niyamas in Modern Life
%%%%%%%%%%%%%%%%%%%%%%%%%%%%%%%%%%%%%%%%%%%%%%%%%%%%%%%%%%%
\begin{frame}[fragile]\frametitle{Yamas \& Niyamas:Ethics for the Digital Age}

\textbf{यम (Yama):Codes of conduct toward others:}
\begin{itemize}
\item \textbf{अहिंसा} Ahimsa:Non-violence: includes online speech, trolling, body shaming
\item \textbf{सत्य} Satya:Truthfulness: versus fake news, performative authenticity on social media
\item \textbf{अस्तेय} Asteya:Non-stealing: time, credit, attention, ideas
\item \textbf{ब्रह्मचर्य} Brahmacharya:Right use of energy: not squandering vitality on endless scrolling
\item \textbf{अपरिग्रह} Aparigraha:Non-possessiveness: versus consumerism, collecting experiences for likes
\end{itemize}

\textbf{नियम (Niyama):Personal observances:}
\begin{itemize}
\item \textbf{शौच} Shaucha:Purity: clean inputs:what you watch, read, listen to
\item \textbf{सन्तोष} Santosha:Contentment: antidote to comparison culture; ``enough'' vs FOMO
\item \textbf{स्वाध्याय} Svadhyaya:Self-study: journaling, therapy, honest self-inquiry
\item \textbf{तपस्} Tapas:Austerity: discipline; doing hard things consistently
\item \textbf{ईश्वरप्रणिधान}:Surrender: releasing the obsessive need to control outcomes
\end{itemize}

\tiny{(Ref: Nancy Wile p.19; Wikipedia:Yoga Sutras; Dr. Ananda Balayogi Bhavanani)}

\end{frame}

%%%%%%%%%%%%%%%%%%%%%%%%%%%%%%%%%%%%%%%%%%%%%%%%%%%%%%%%%%%
%% NEW SLIDE C4: Breaking the Karma Loop
%%%%%%%%%%%%%%%%%%%%%%%%%%%%%%%%%%%%%%%%%%%%%%%%%%%%%%%%%%%
\begin{frame}[fragile]\frametitle{2.12–2.17:Breaking the Karma Loop}

\begin{sanskrit}
हेयं दुःखमनागतम् ॥२.१६॥
द्रष्टृदृश्ययोः संयोगो हेयहेतुः ॥२.१७॥
\end{sanskrit}

\begin{itemize}
\item \textbf{2.12:} Latent impressions (samskaras) colored by kleshas ripen into experience:in this life or future ones. All our actions are stored as seeds in the mind.
\item \textbf{2.13–2.14:} As long as kleshas remain at the root → birth, span of life, and experiences follow as pleasure or pain
\item \textbf{2.17:The root cause:} The alliance between the Seer (subject) and the Seen (object) is the cause of suffering. Yoga breaks this alliance.
\end{itemize}

\textbf{You always have a choice (Nancy Wile):}
\begin{itemize}
\item Car breaks down in rain: \emph{play the victim} (anger + suffering) OR \emph{accept and solve} (same outcome, different inner experience)
\item ``Two people can have the same thought but very different reactions. Why?'':because of their samskaras
\item Good thoughts and kind actions eventually return to us; fearful thoughts attract what we fear
\item \textbf{Yoga = planting new seeds in the mind for a world free of pain} (Nancy Wile, 2.18)
\end{itemize}

\tiny{(Ref: Nancy Wile p.16-17; Dr. Ananda Balayogi Bhavanani; [IT] I.K. Taimni)}

\end{frame}

%%%%%%%%%%%%%%%%%%%%%%%%%%%%%%%%%%%%%%%%%%%%%%%%%%%%%%%%%%%
\begin{frame}[fragile]\frametitle{2.28–2.29:Ashtanga Yoga: The Eight-Limbed Path}

\begin{sanskrit}
योगाङ्गानुष्ठानादशुद्धिक्षये ज्ञानदीप्तिराविवेकख्यातेः ॥२.२८॥
यमनियमासनप्राणायामप्रत्याहारधारणाध्यानसमाधयोऽष्टावङ्गानि ॥२.२९॥
\end{sanskrit}

\textbf{2.28:} Sustained practice of all 8 limbs destroys impurities and cultivates the light of wisdom leading to Viveka-khyati (discriminative knowledge).

\textbf{The Eight Limbs (अष्टांग):}
\begin{center}
{\small
\begin{tabular}{|p{0.5cm}|p{2.5cm}|p{4cm}|}
\hline
\textbf{\#} & \textbf{Limb} & \textbf{Domain} \\
\hline
1 & यम Yama (5) & Restraints: Ahimsa, Satya, Asteya, Brahmacharya, Aparigraha \\
2 & नियम Niyama (5) & Observances: Shaucha, Santosha, Tapas, Svadhyaya, Ishvara-pranidhana \\
3 & आसन Asana & Steady, comfortable posture \\
4 & प्राणायाम Pranayama & Breath regulation \\
5 & प्रत्याहार Pratyahara & Withdrawal of senses \\
6 & धारणा Dharana & Concentration \\
7 & ध्यान Dhyana & Meditation \\
8 & समाधि Samadhi & Absorption \\
\hline
\end{tabular}
}
\end{center}

\tiny{(Ref: Nancy Wile; Dr. Ananda Balayogi Bhavanani; [HA] Hariharananda Aranya)}

\end{frame}

%%%%%%%%%%%%%%%%%%%%%%%%%%%%%%%%%%%%%%%%%%%%%%%%%%%%%%%%%%%
\begin{frame}[fragile]\frametitle{2.46, 2.49, 2.54:Outer Limbs: Asana, Pranayama, Pratyahara}

\begin{sanskrit}
स्थिरसुखमासनम् ॥२.४६॥
\end{sanskrit}

\textbf{2.46:Asana:} ``Posture is that which is \textbf{steady and comfortable}.'' Patanjali lists no specific forms:the key is \emph{relaxation of effort} and \emph{contemplation of the infinite} (2.47). Through asana one transcends the dualities that normally torment us (2.48).

\begin{sanskrit}
तस्मिन् सति श्वासप्रश्वासयोर्गतिविच्छेदः प्राणायामः ॥२.४९॥
\end{sanskrit}

\textbf{2.49:Pranayama:} Cessation of inhalation and exhalation. Dr. Ananda: ``When amazed, we say \emph{It took my breath away!} Imagine the state of breath face-to-face with the Divine.''

\textbf{2.54:Pratyahara:} Withdrawal of senses from their objects. Like a tortoise pulling its limbs into its shell. ``Not consciously closing eyes to the sensory world:consciously closing the \emph{mind processes} to the sensory world.'' (Wikipedia) This is the bridge from outer (bahiranga) to inner (antaranga) Yoga.

\tiny{(Ref: Dr. Ananda Balayogi Bhavanani; Nancy Wile; Wikipedia; [SS] Swami Satchidananda)}

\end{frame}

%%%%%%%%%%%%%%%%%%%%%%%%%%%%%%%%%%%%%%%%%%%%%%%%%%%%%%%%%%%%%%%%%%%%%%%%%%%%%%%%%%
%% SECTION: Chapter 3:Vibhuti Pada
%%%%%%%%%%%%%%%%%%%%%%%%%%%%%%%%%%%%%%%%%%%%%%%%%%%%%%%%%%%%%%%%%%%%%%%%%%%%%%%%%%
\section[Vibhutipad]{Chapter 3: Vibhuti Pada:Experiences}

\begin{frame}[fragile]\frametitle{Chapter 3: विभूतिपाद Vibhuti Pada:Experiences \& Powers (56 Sutras)}

\begin{itemize}
\item \textbf{Theme:} The inner (antaranga) limbs:Dharana, Dhyana, Samadhi:and the extraordinary fruits of their combined practice (Samyama)
\end{itemize}

\begin{sanskrit}
देशबन्धश्चित्तस्य धारणा ॥३.१॥
तत्र प्रत्ययैकतानता ध्यानम् ॥३.२॥
तदेवार्थमात्रनिर्भासं स्वरूपशून्यमिव समाधिः ॥३.३॥
तत्रयमेकसंयमः ॥३.४॥
\end{sanskrit}

\begin{itemize}
\item \textbf{3.1:धारणा:} binding consciousness to a point, place, or object
\item \textbf{3.2:ध्यान:} steady, continuous flow of attention:Dharana sustained without break
\item \textbf{3.3:समाधि:} object alone shines; meditator loses separate identity; ``Thinker, thought, and thing thought about fuse into one.'' (Wikipedia)
\item \textbf{3.4:संयम:} all three practiced together as one seamless movement:the master tool of this chapter
\end{itemize}

\tiny{(Ref: Dr. Ananda Balayogi Bhavanani; [VH] Vyasa Houston; Wikipedia)}

\end{frame}

%%%%%%%%%%%%%%%%%%%%%%%%%%%%%%%%%%%%%%%%%%%%%%%%%%%%%%%%%%%
%% NEW SLIDE D1: Samyama as Flow State
%%%%%%%%%%%%%%%%%%%%%%%%%%%%%%%%%%%%%%%%%%%%%%%%%%%%%%%%%%%
\begin{frame}[fragile]\frametitle{Samyama = Flow State:Ancient Name for a Modern Experience}

\begin{columns}
\begin{column}[t]{0.5\linewidth}
\textbf{Patanjali's Samyama (3.4)}
\begin{itemize}
\item Dharana: fix attention on one object
\item Dhyana: sustained, unbroken flow
\item Samadhi: merge with the object; self disappears
\item Result: direct knowledge of the object's essence
\end{itemize}
\end{column}
\begin{column}[t]{0.5\linewidth}
\textbf{Csikszentmihalyi's Flow}
\begin{itemize}
\item Total absorption in one activity
\item Effortless attention; time distorts
\item Sense of self drops away
\item Result: peak performance + deep satisfaction
\end{itemize}
\end{column}
\end{columns}

\vspace{3mm}
\begin{itemize}
\item \textbf{You already know Samyama:} when coding, playing music, painting, or playing sport:and hours pass like minutes:that is Samyama in action
\item The difference: Yoga makes this entry into absorption \emph{intentional, reproducible, and direction-able} toward the subtlest levels of reality
\item \textbf{3.5–3.6:} Mastery of Samyama brings light of knowledge; progressively applied to subtler planes
\end{itemize}

\tiny{(Ref: Mihaly Csikszentmihalyi:\emph{Flow}; Dr. Ananda Balayogi Bhavanani; Wikipedia)}

\end{frame}

%%%%%%%%%%%%%%%%%%%%%%%%%%%%%%%%%%%%%%%%%%%%%%%%%%%%%%%%%%%
\begin{frame}[fragile]\frametitle{Siddhis:Attainments and Obstacles (3.38, 3.50–3.51)}

\begin{itemize}
\item \textbf{3.17–3.37 and 3.39–3.49:} Samyama on various objects yields extraordinary powers (siddhis/vibhutis): knowledge of past lives, reading minds, levitation, mastery over elements, etc.
\end{itemize}

\begin{sanskrit}
ते समाधावुपसर्गा व्युत्थाने सिद्धयः ॥३.३८॥
\end{sanskrit}

\begin{itemize}
\item \textbf{3.38:Critical Warning:} Siddhis are \textbf{attainments} in the waking state but \textbf{obstacles} to Samadhi. ``The temptation of these powers should be avoided and the attention fixed only on liberation.'':Nancy Wile
\end{itemize}

\begin{sanskrit}
तद्वैराग्यादपि दोषबीजक्षये कैवल्यम् ॥३.५१॥
\end{sanskrit}

\begin{itemize}
\item \textbf{3.50–3.51:Para-Vairagya:} Kaivalya is attained only through renunciation even of the highest siddhis. Decline even celestial invitations:anything less than Kaivalya is still bondage.
\item \textbf{3.56:} Liberation comes when Buddhi and Purusha become equal in purity:total discrimination achieved.
\end{itemize}

\tiny{(Ref: Dr. Ananda Balayogi Bhavanani; Nancy Wile; yogasutrastudy.info)}

\end{frame}

%%%%%%%%%%%%%%%%%%%%%%%%%%%%%%%%%%%%%%%%%%%%%%%%%%%%%%%%%%%%%%%%%%%%%%%%%%%%%%%%%%
%% SECTION: Chapter 4:Kaivalya Pada
%%%%%%%%%%%%%%%%%%%%%%%%%%%%%%%%%%%%%%%%%%%%%%%%%%%%%%%%%%%%%%%%%%%%%%%%%%%%%%%%%%
\section[Kaivalyapad]{Chapter 4: Kaivalya Pada:Absolute Freedom}

\begin{frame}[fragile]\frametitle{Chapter 4: कैवल्यपाद Kaivalya Pada:Absolute Freedom (34 Sutras)}

\begin{itemize}
\item \textbf{Theme:} The nature of consciousness, liberation, and what lies beyond all experience
\item Kaivalya = ``aloneness'':not isolation, but the deepest realization where no division between self and other exists
\item ``What happens when you take control of your thoughts rather than having your thoughts control you.'':Nancy Wile
\end{itemize}

\begin{sanskrit}
जन्मौषधिमन्त्रतपःसमाधिजाःसिद्धयः ॥४.१॥
\end{sanskrit}

\begin{itemize}
\item \textbf{4.1:} Siddhis may arise from birth, herbs, mantras, austerities, or Samadhi:but they are never the goal
\item \textbf{4.4–4.6:} All mind-constructs ultimately flow from the ego (asmita); the liberated yogi is free from ego-driven mind
\item \textbf{4.15–4.17:} The same object is perceived differently by different minds:``No wonder everyone seems to have his or her own view of the world.'' (Dr. Ananda)
\end{itemize}

\tiny{(Ref: Dr. Ananda Balayogi Bhavanani; Nancy Wile; Wikipedia)}

\end{frame}

%%%%%%%%%%%%%%%%%%%%%%%%%%%%%%%%%%%%%%%%%%%%%%%%%%%%%%%%%%%
\begin{frame}[fragile]\frametitle{4.7–4.12:Karma and Samskaras}

\begin{sanskrit}
कर्माशुक्लाकृष्णं योगिनस्त्रिविधमितरेषाम् ॥४.७॥
जातिदेशकालव्यवहितानामप्यानन्तर्यं स्मृतिसंस्कारयोरेकरूपत्वात् ॥४.९॥
\end{sanskrit}

\begin{itemize}
\item \textbf{4.7:Types of karma:} For ordinary people, karma is white/black/mixed. For the yogi, karma is neither:they act without binding consequences.
\item \textbf{4.8:} Karma ripens only when the right conditions of species, place, and time arise.
\item \textbf{4.9:Samskaras across lifetimes:} Memory and latent impressions share one form:unbroken continuity \emph{across births}. ``Deep habit patterns play out from lifetime to lifetime.'' (Dr. Ananda)
\item \textbf{4.10:} The will to live (Abhinivesha) is beginningless:confirming the reality of the cycle.
\item \textbf{4.12:} Past and future both exist in the present, differing only in characteristic:by knowing the present fully, one may know past and future.
\item \textbf{4.27:} Even at the highest level, deep samskaras can re-surface and stall progress:remain vigilant (Dr. Ananda).
\end{itemize}

\tiny{(Ref: Dr. Ananda Balayogi Bhavanani; [IT] I.K. Taimni; Nancy Wile)}

\end{frame}

%%%%%%%%%%%%%%%%%%%%%%%%%%%%%%%%%%%%%%%%%%%%%%%%%%%%%%%%%%%
\begin{frame}[fragile]\frametitle{4.29–4.34:Dharma Megha Samadhi and Kaivalya}

\begin{sanskrit}
प्रसंख्यानेऽप्यकुसीदस्य सर्वथा विवेकख्यातेर्धर्ममेघः समाधिः ॥४.२९॥
ततः क्लेशकर्मनिवृत्तिः ॥४.३०॥
\end{sanskrit}

\textbf{4.29:Dharma Megha Samadhi (धर्ममेघसमाधि):}
\begin{itemize}
\item ``The rain cloud of virtue'':even attachment to discriminative knowledge is released; the mind pours forth righteousness naturally
\item ``The torrential rainfall from this rain cloud of the highest nature washes away all arrogant impurities.'':Dr. Ananda
\end{itemize}

\textbf{4.30:} All kleshas and karma removed once and for all:the final liberation. ``At this point, Patanjali implies we become the Divine itself in the state of Kaivalya.'' (Dr. Ananda)

\begin{sanskrit}
पुरुषार्थशून्यानां गुणानां प्रतिप्रसवः कैवल्यं स्वरूपप्रतिष्ठा वा चितिशक्तिरिति ॥४.३४॥
\end{sanskrit}

\textbf{4.34:Kaivalya:} The gunas resolve back into their cause. Pure consciousness is established in its own true nature. The evolutionary journey ends. \textbf{Sat-Chit-Ananda}. ``We lose our sense of individuality in order to gain the sense of absolute universality.'' (Dr. Ananda)

\tiny{(Ref: Dr. Ananda Balayogi Bhavanani; [HA] Hariharananda Aranya; [IT] I.K. Taimni)}

\end{frame}

%%%%%%%%%%%%%%%%%%%%%%%%%%%%%%%%%%%%%%%%%%%%%%%%%%%%%%%%%%%%%%%%%%%%%%%%%%%%%%%%%%
%% SECTION: Conclusion
%%%%%%%%%%%%%%%%%%%%%%%%%%%%%%%%%%%%%%%%%%%%%%%%%%%%%%%%%%%%%%%%%%%%%%%%%%%%%%%%%%
\section[Conclusion]{The Complete Journey}

%%%%%%%%%%%%%%%%%%%%%%%%%%%%%%%%%%%%%%%%%%%%%%%%%%%%%%%%%%%
%% NEW SLIDE E1: Cross-tradition Philosophical Context
%%%%%%%%%%%%%%%%%%%%%%%%%%%%%%%%%%%%%%%%%%%%%%%%%%%%%%%%%%%
\begin{frame}[fragile]\frametitle{Philosophical Roots:A Universal Science}

\begin{itemize}
\item \textbf{Samkhya-Yoga:} Samkhya = \emph{theory} (why we suffer); Yoga = \emph{practice} (how to be free). The Bhagavad Gita is built on this Samkhya-Yoga foundation.
\item \textbf{Buddhism parallel:} Ashtanga Yoga's eight limbs parallel the Buddhist \emph{Noble Eightfold Path}. Brahmaviharas (1.33: Maitri, Karuna, Mudita, Upeksha) appear in both traditions:a universal human wisdom.
\item \textbf{Jainism parallel:} The five Yamas closely resemble the five major vows of Jainism. Ahimsa (non-violence) appears first in both.
\item \textbf{Modern psychology:} Klesha theory parallels cognitive distortions in CBT; Pratipaksha Bhavana = reframing; the five Chitta states = attention research.
\end{itemize}

\begin{center}
{\small ``The Yoga Sutras clarify esoteric concepts that are common to \emph{all traditions of Indian thought}.'' (Nancy Wile)\\
\textbf{Yoga does not belong to any religion. It belongs to humanity.}}
\end{center}

\tiny{(Ref: Wikipedia:Philosophical Roots; Nancy Wile; Dr. Ananda Balayogi Bhavanani)}

\end{frame}

%%%%%%%%%%%%%%%%%%%%%%%%%%%%%%%%%%%%%%%%%%%%%%%%%%%%%%%%%%%
\begin{frame}[fragile]\frametitle{The Complete Journey: Vritti to Kaivalya}

\begin{center}
\begin{tabular}{c@{\ $\longrightarrow$\ }l}
\textbf{Vritti} & Mind in constant modification \\
\textbf{Coloring} & Thoughts colored by kleshas and past samskaras \\
\textbf{Nirodha} & Restraint through Abhyasa + Vairagya \\
\textbf{Samadhi} & Deepening absorption; crystal mind (Ch.1) \\
\textbf{Kriya Yoga} & Purifying afflictions; breaking karma loop (Ch.2) \\
\textbf{Ashtanga} & 8-limbed discipline; digital ethics and inner life (Ch.2) \\
\textbf{Samyama} & Inner mastery; flow state intentionalized (Ch.3) \\
\textbf{Viveka-khyati} & Discriminative knowledge (Ch.3–4) \\
\textbf{Dharma Megha} & Rain cloud of virtue (Ch.4) \\
\textbf{Kaivalya} & Purusha in its own nature:absolute freedom (Ch.4) \\
\end{tabular}
\end{center}

\vspace{3mm}
\begin{center}
{\large ``Yoga begins when you become ready to be taught.\\It ends when you realize you were never what you thought you were.''}
\end{center}

\end{frame}

%%%%%%%%%%%%%%%%%%%%%%%%%%%%%%%%%%%%%%%%%%%%%%%%%%%%%%%%%%%
%% NEW SLIDE F1: Start Today
%%%%%%%%%%%%%%%%%%%%%%%%%%%%%%%%%%%%%%%%%%%%%%%%%%%%%%%%%%%
\begin{frame}[fragile]\frametitle{Start Today:Your Personal Sadhana}

\begin{center}
{\normalsize \textbf{You do not need to master 196 sutras to begin. Start with one.}}
\end{center}

\vspace{3mm}
\textbf{A simple daily practice (10–15 min):}
\begin{enumerate}
\item \textbf{Morning:} Choose one method from 1.34–1.39 that feels natural:pranayama, inner light, or contemplating what you love most
\item \textbf{One sutra a day:} Read one sutra, sit quietly with it before moving to the next. ``Get into a quiet posture and notice what meaning you get from each one.'':Nancy Wile
\item \textbf{In relationships:} Practice Maitri-Karuna-Mudita-Upeksha in one interaction daily (1.33)
\item \textbf{When negative thought arises:} Pratipaksha Bhavana:say STOP, then cultivate the opposite (2.33)
\item \textbf{Evening:} Ask: ``Were my thoughts today facts or colored concepts?'' (1.5-1.6)
\end{enumerate}

\vspace{2mm}
\begin{center}
{\small ``When that practice is done for a long time, without a break, and with sincere devotion, then it becomes a firmly rooted, stable and solid foundation.'':1.14}
\end{center}

\tiny{(Ref: Nancy Wile, Yoga Education Institute, p.4; Patanjali 1.14)}

\end{frame}

%%%%%%%%%%%%%%%%%%%%%%%%%%%%%%%%%%%%%%%%%%%%%%%%%%%%%%%%%%%
%% NEW SLIDE F2: 10 Sutras to Carry
%%%%%%%%%%%%%%%%%%%%%%%%%%%%%%%%%%%%%%%%%%%%%%%%%%%%%%%%%%%
\begin{frame}[fragile]\frametitle{10 Sutras to Carry:Pocket Wisdom for Daily Life}

\begin{center}
{\small
\begin{tabular}{|p{1cm}|p{1.5cm}|p{4.5cm}|}
\hline
\textbf{Sutra} & \textbf{Key word} & \textbf{What it says to you today} \\
\hline
1.1 & अथ & You are ready. Begin now. \\
1.2 & निरोधः & You are not your thoughts. \\
1.12 & अभ्यास-वैराग्य & Practice and let go; both wings needed. \\
1.14 & दृढभूमिः & Long, unbroken, sincere effort builds everything. \\
1.33 & मैत्री & Friendliness, compassion, joy, equanimity; practice today. \\
1.39 & यदभिमतम् & Whatever you love most can be your path to stillness. \\
2.16 & अनागतम् & The pain yet to come can be avoided; act now. \\
2.33 & प्रतिपक्ष & When a negative thought grips you, cultivate the opposite. \\
3.38 & उपसर्गाः & Any attainment that distracts from liberation is an obstacle. \\
4.34 & कैवल्यम् & The evolutionary journey ends in absolute freedom. \\
\hline
\end{tabular}
}
\end{center}

\tiny{(Ref: Compiled from Nancy Wile; Dr. Ananda Balayogi Bhavanani; Krishnamurthy Ramakrishnan)}

\end{frame}

%%%%%%%%%%%%%%%%%%%%%%%%%%%%%%%%%%%%%%%%%%%%%%%%%%%%%%%%%%%
%%%%%%%%%%%%%%%%%%%%%%%%%%%%%%%%%%%%%%%%%%%%%%%%%%%%%%%%%%%%%%%%%%%%%%%%%%%%%%%%%%
\begin{frame}[fragile]\frametitle{}
\begin{center}
{\Large References}
\end{center}
\end{frame}


%%%%%%%%%%%%%%%%%%%%%%%%%%%%%%%%%%%%%%%%%%%%%%%%%%%%%%%%%%%
\begin{frame}[fragile]\frametitle{References}

Many publicly available sources have been used in the preparation of this content. Some of the salient ones are listed below:

	\begin{itemize}
	\item Patanjali Yog Sutra https://patanjaliyogasutra.in/
	\item Yoga Sutra Study https://yogasutrastudy.info/
		\begin{itemize}
		\item [HA]: Hariharananda Aranya
		\item [IT]: I. K. Taimni
		\item [VH]: Vyasa Houston
		\item [BM]: Barbara Miller
		\item [SS]: Swami Satchidananda
		\item [SP]: Swami Prabhavananda
		\item [SV]: Swami Vivekananda
		\end{itemize}
	\item Prof. Edwin Bryant — \emph{The Yoga Sutras of Patanjali} (book + lectures)
	\item Dr. Ananda Balayogi Bhavanani — \emph{An Overview of the Yoga Sutras}, Integral Yoga Magazine
	\item Nancy Wile — \emph{The Yoga Sutras} (Yoga Teacher Training manual), Yoga Education Institute, 2014
	\item Krishnamurthy Ramakrishnan — \emph{Yoga Sutras for Daily Life}, LinkedIn Newsletter
	\item Wikipedia — \emph{Yoga Sutras of Patanjali}
	\item योग सूत्र का परिचय | Yoga Sutras of Patanjali | Yoga Sutras For QCI, UGC NET, Yoga All Exams - YouTube
	\item Patanjali Yoga Sutra Dr Mrudula Chaudhari
	\item Patanjali Yoga Sutras | Explanation by Anandmurti Gurumaa Gurumaa Ashram Verified
	\item Indica Course https://indica.courses/enroll/prof-m-a-narasimhan/maharshi-patanjali-yogasutras-learn-recite-ii/
	\item Swami J Patanjali Yoga Sutra https://swamij.com/yoga-sutras.htm
	\item Dr. Veena Londhe — \emph{An Introduction to Patanjali's Yoga Darshana}, Chinmaya Mission
	\end{itemize}

\end{frame}

%%%%%%%%%%%%%%%%%%%%%%%%%%%%%%%%%%%%%%%%%%%%%%%%%%%%%%%%%%%
\begin{frame}[fragile]\frametitle{Youtube Channels}



	\begin{itemize}
	\item Patanjali Yoga Sutras - English Yogic Gurukul - Anvita Dixit · Podcast
	\item Patanjali Yoga Sutras | Explanation by Anandmurti Gurumaa
	\item Complete Patanjali Yoga Sutras Chant with Meanings
	\item Patanjali Yoga Sutra Live Streams - The Sanskrit Channel
	\item The Sankrit Channel - "Detailed Patanjali Yoga Sutras with Pictures - Chapter 1-2-3-4"
	\end{itemize}

\end{frame}

%%%%%%%%%%%%%%%%%%%%%%%%%%%%%%%%%%%%%%%%%%%%%%%%%%%%%%%%%%%
\begin{frame}[fragile]\frametitle{}

\begin{center}
\includegraphics[width=0.8\linewidth,keepaspectratio]{my_yog_back}

Thanks धन्यवाद
\end{center}

\end{frame}

